\chapter{Visualización del movimiento en mapas interactivos}
\label{ch:o5}

El capítulo anterior entregó, para cada par (ave, día) del conjunto de prueba, una
predicción continua del movimiento del día siguiente, con un punto central y una banda
de incertidumbre por componente. Este capítulo describe la herramienta que lleva esa
predicción al mapa: un módulo de visualización que representa sobre el territorio las
rutas de las aves y su predicción, y que permite explorar de forma interactiva, en
kilómetros y de un vistazo, el error de cada predicción. La herramienta no entrena
ningún modelo nuevo ni reabre el análisis cuantitativo del error, ya realizado en el
capítulo \ref{ch:o4}; su aportación es hacer ese comportamiento tangible y explorable,
y con ella se cierra la cadena que recorre la memoria, de la cadena de Markov visible al
modelo oculto y de este a la predicción supervisada. Las secciones siguientes describen
qué ofrece la herramienta y cómo se usa, y después, brevemente, cómo se ha construido.

\section{Funcionalidades y uso}
\label{sec:o5-funcionalidades}

La herramienta es una aplicación web interactiva (figura \ref{fig:o5-app}) que se abre
directamente en el navegador, sin instalación ni servidor, y que representa sobre el mapa
el recorrido real y la predicción del modelo para cualquiera de las ochenta y dos aves
del conjunto de prueba. Toma como base la predicción recomendada del capítulo
\ref{ch:o4} y reúne en un único panel lateral las funcionalidades de exploración que se
resumen a continuación:

\begin{itemize}
  \item \textbf{Selección de ave.} Cualquiera de las aves del conjunto, ordenadas por
  histórico, más una vista «Todas» que dibuja a la vez las rutas reales de todas ellas
  (un color por ave, con resalte al pasar el ratón y acceso a la vista individual al
  pulsar sobre una de las rutas).
  \item \textbf{Predicción del día y error en kilómetros.} Para el ave elegida se
  superponen su posición real, el punto predicho a un día y su banda de incertidumbre
  $[p_{10}, p_{90}]$; una lectura destacada muestra el error de ese día en kilómetros,
  esto es, la distancia entre la predicción y la posición real. Esta lectura cumple
  directamente el cometido de analizar visualmente el error de predicción.
  \item \textbf{Reproducción temporal.} Una ventana de días (con inicio y fin
  ajustables) y los controles habituales de reproducir, pausar, avanzar o retroceder un
  día y regular la velocidad; al reproducir, la estela de la predicción va creciendo
  sobre el mapa.
  \item \textbf{Capas conmutables.} La trayectoria real completa del ave (coloreada por
  año), el punto del baseline de Markov de primer orden con su propio error, y el estado
  de comportamiento del modelo oculto del capítulo \ref{ch:o3} (estacionario o
  migración).
\end{itemize}

El uso típico es inmediato: se abre el fichero en el navegador, se elige un ave, se
ajusta la ventana de días y se pulsa reproducir para seguir, día a día, cómo la
predicción persigue a la posición real, leyendo en todo momento el error en kilómetros.
Para que esa lectura no sea ambigua, la herramienta fija un código de color estable, con
el error siempre en rojo, la posición real en verde y la predicción en granate. La misma
herramienta incluye, además, una demostración de predicción encadenada a varios pasos,
etiquetada como ilustrativa por las razones de calibración que se explican más adelante.

\begin{figure}[H]
  \centering
  \includegraphics[width=\linewidth]{o5_prediccion_app.png}
  \caption[Aplicación interactiva de visualización de la predicción]{Captura de la
  aplicación interactiva a medida. El panel lateral reúne la selección de ave, la
  ventana y los controles de reproducción y los conmutadores de capas; el mapa muestra,
  para el día actual, la posición real, la predicción a un día con su banda de
  incertidumbre y la lectura del error en kilómetros. El fichero es una página web
  autónoma, de la que aquí se reproduce una captura estática.}
  \label{fig:o5-app}
\end{figure}

\section{Desarrollo}
\label{sec:o5-desarrollo}

El módulo se implementa en el paquete \texttt{tfg\_aves.viz}, con el mismo patrón que el
resto del trabajo: el cálculo de los agregados, en módulos de cómputo puro, se separa
del renderizado. La aplicación es una plantilla HTML que contiene todo el código de la
interfaz (estilos y lógica de cliente, sobre la biblioteca de mapas Leaflet) y en la que
se inyecta, como datos embebidos en formato JSON, la predicción de cada ave; que sea una
página estática sin servidor es posible porque las predicciones del conjunto de prueba
están precalculadas.

Sobre qué datos se dibujan se toman tres decisiones técnicas. La herramienta cartografía
el predictor recomendado por el capítulo \ref{ch:o4} (la regresión de cuantiles con
LightGBM en modo poblacional). La incertidumbre se dibuja como
el rectángulo $[p_{10}, p_{90}]$ de cada eje, y no como una elipse, porque es la
representación fiel de unos cuantiles estimados por componente. Y el horizonte de la
vista es de un solo día, el único en el que la banda está calibrada; el encadenamiento a
varios pasos se ofrece únicamente como la demostración ilustrativa ya mencionada.

% ===========================================================================
% PENDIENTE (no redactar hasta aprobación de las secciones anteriores):
%   \section{Síntesis ...}   -> cierre breve: qué hace evidente la herramienta,
%                               límites, enganche con conclusiones.
% ===========================================================================
